\subsection{不变双线性型}

命题 5. 设 \((\mathfrak{g},\mathfrak{h})\) 是分裂半单李代数，\(\Phi\) 是 g 上的不变对称双线性型，W 是 \((\mathfrak{g},\mathfrak{h})\) 的 Weyl 群。则 \(\Phi'\) （\(\Phi\) 在 h 上的限制）在 W 下不变。而且，若 \(\Phi\) 非退化，则 \(\Phi'\) 也非退化。

设 \(\alpha \in \mathbb{R}\) ，设 \(X_{\alpha}\) 是 \(\mathfrak{g}^{\alpha}\) 的非零元素，\(\rho\) 是对应的 \(\mathfrak{sl}(2,k)\) 在 \(\mathfrak{g}\) 上的表示，\(\pi\) 是 \(\mathbf{SL}(2,k)\) 在 \(\mathfrak{g}\) 上与 \(\rho\) 相容的表示。则 \(\varPhi\) 在 \(\rho\) 下不变，从而在 \(\pi\) 下不变（§1 第 4 节）。特别地，\(\varPhi'\) 在 \(\theta_{\alpha}(t)|\mathfrak{h}\) （第 2 节）下不变，从而在 W 下不变。最后一个断言由命题 1 (i) 得出。

命题 6. 设 \((\mathfrak{g},\mathfrak{h})\) 是分裂半单李代数，\(\varPhi\) 是 \(\mathfrak{g}\) 上的非退化不变对称双线性型。对一切 \(\alpha \in \mathbb{R}\) ，设 \(X_{\alpha}\) 是 \(\mathfrak{g}^{\alpha}\) 的非零元素。设 \((H_i)_{i\in I}\) 是 \(\mathfrak{h}\) 的基，\((H_i')_{i\in I}\) 是 \(\mathfrak{h}\) 中使 \(\varPhi(H_i,H_j') = \delta_{ij}\) 的基。这时，对应于 \(\varPhi\) 的 Casimir 元（\(\mathfrak{g}\) 的包络代数中的，第一章 §3 第 7 节）是

\[
\sum_ {\alpha \in \mathrm{R}} \frac {1}{\varPhi (X _ {\alpha} , X _ {- \alpha})} X _ {\alpha} X _ {- \alpha} + \sum_ {i \in I} H _ {i} H _ {i} ^ {\prime}.
\]

事实上，由命题 1，\(\Phi(H_i, X_\alpha) = \Phi(H_i', X_\alpha) = 0\) 对一切 \(i \in I\) ，\(\alpha \in \mathrm{R}\) 成立，并且 \(\Phi\left(\frac{1}{\Phi(X_\alpha, X_{-\alpha})} X_\alpha, X_{-\beta}\right) = \delta_{\alpha\beta}\) 对一切 \(\alpha, \beta \in \mathrm{R}\) 成立。

